\subsubsection*{Remark 5.1.}
Condition

\begin{equation}\tag{5.1.1}
jd<n+2
\end{equation}

admits a cohomological interpretation.  Let $V$ be a smooth hypersurface of
degree $d$ and dimension $n$.  As in XI, 2.1, put, for $a+b=n$,

\begin{equation}\tag{5.1.2}
h_0^{a,b}
=\dim H^b(V,\Omega_V^a)-\delta_{a,b}
=\dim E^{a,b}(V),
\end{equation}

where $E^{a,b}(V)$ is the bidegree-$(a,b)$ component of the vanishing Hodge
cohomology (compare XVII, 5.4.4)

\begin{equation}\tag{5.1.3}
E^{a,b}(V)
=\operatorname{Coker}\!\left(
H^b(\mathbb P^{n+1},\Omega_{\mathbb P^{n+1}}^a)
\longrightarrow H^b(V,\Omega_V^a)\right).
\end{equation}

Then, by XI, 2.8,

\begin{equation}\tag{5.1.4}
jd<n+2
\quad\Longleftrightarrow\quad
h_0^{a,n-a}=0\quad\text{for }a<j\text{ and }a>n-j.
\end{equation}

In other words, the condition says that the Hodge level (XI, 2.7) of the
vanishing de Rham cohomology of $V$,

\[
E_{\mathrm{DR}}^n(V)
=\operatorname{Coker}\!\left(
H_{\mathrm{Hdg}}^n(\mathbb P^{n+1})
\longrightarrow H_{\mathrm{Hdg}}^n(V)\right),
\]

is at most $n-2j$, or equivalently that its coniveau is at least $j$.
On the other hand, the conclusion also says that the eigenvalues of the
Frobenius $\mathfrak F_q$ acting on the vanishing $\ell$-adic cohomology
($\ell\ne p$) are algebraic integers divisible by $q^j$ (XXI, A.2).
Thus Ax's result 5.0 supports the following conjecture, suggested by the
philosophy of the level of motives and relating the classical Hodge and
Tate conjectures~[5].

\subsubsection*{Conjecture 5.2.}
Let $f:X\to S$ be a proper smooth morphism, with $S$ integral, of finite
type over $\operatorname{Spec}(\mathbb Z)$, and dominant over
$\operatorname{Spec}(\mathbb Z)$; let $i,j$ be integers.  Suppose that the
Hodge coniveau of $H^i$ of the generic fiber of $f$ is at least $j$.  Then,
for every point $s\in S$, the coniveau (XX, 2.0.6) of
$H^i(X_s,\mathbb Q_\ell)$ is at least $j$.  Hence (XX, 2.2), if $s$ is a
closed point of $S$, the eigenvalues of Frobenius acting on
$H^i(X_s,\mathbb Q_\ell)$ are algebraic integers divisible by $q^j$, where
$q=\operatorname{card}(k(s))$.

\section*{6. End of the proof of the congruence formula}

Taking account of (4.6.6) and (4.6.3)--(4.6.4), formula (4.5.33) gives

\begin{equation}\tag{6.0.1}
Z(t,V/\mathbb F_q)(1-t)
\equiv\Delta_S(t)^{(-1)^{n+1}}
\pmod{1+qt\mathcal O[[t]]}.
\end{equation}

We now calculate $\Delta_S(t)$ modulo $\pi\mathcal O[[t]]$.  Define, as in
(4.6.11),

\begin{equation}\tag{6.0.2}
W_S=R_S\otimes_{\mathcal O}\mathbb F_q.
\end{equation}

\medskip
\noindent\emph{Source note.}
The print cites (4.5.24), but the congruence follows from the terminal zeta
identity (4.5.33).  It also defines $W_S$ using $R_B$, although the entire
section and every following formula concern $R_S$.  The corrected French
and English layers use (4.5.33) and $R_S$; the diplomatic French retains
the printed forms.
